\section{Teorema rector de la red radical y la ley potencia--vacancia}

La aparición de \(369\) y \(693\) no se reduce a una semejanza gráfica. Hay dos mecanismos exactos que convergen:

\begin{enumerate}
\item \textbf{mecanismo algebraico nonádico:} \(3,6,9\) son los representantes visuales del radical nilpotente \(\mathfrak m=(3)\subset\mathbb Z/9\mathbb Z\); las palabras \(369\), \(693\) y \(936\) son las tres fases de su órbita multiplicativa;
\item \textbf{mecanismo arquimediano de vacancias:} la diferencia entre la longitud decimal de \(9^n\) y la frontera \(10^n\) está gobernada exactamente por la misma rotación irracional \(\delta_{\rm vac}=\log_{10}(10/9)\) que ya organiza el calendario de vacancias.
\end{enumerate}

Para \(n=9^9\), ambos mecanismos se encuentran en la identidad

\[
9^9-\left\lceil 9^9\log_{10}\frac{10}{9}\right\rceil
=369\,693\,099,
\]

de modo que

\[
\#\operatorname{cifras}\bigl(9^{9^9}\bigr)=369\,693\,100.
\]

El bloque \(369\mid693\) es, por tanto, una salida exacta del lector logarítmico de la potencia nonádica; no es el prefijo decimal de \(9^{9^9}\). La unidad final procede del operador universal que transforma orden decimal en longitud decimal.

Esta identidad abre una lectura estructural de mayor alcance: suma, producto, potencia, resta, división, raíz y logaritmo pueden organizarse como una jerarquía de transportes y cocientes. Cada ascenso de operación compactifica historias; el TPK conserva las coordenadas que permiten reconstruirlas.


\section{Estructura radical de la APP sobre \(\mathbb Z/9\mathbb Z\)}

\subsection{Anillo local nonádico \(\mathbb Z/9\mathbb Z\) y radical nilpotente}

Sea

\[
R=\mathbb Z/9\mathbb Z.
\]

Su grupo de unidades y su ideal maximal son

\[
R^\times=\{1,2,4,5,7,8\},
\qquad
\mathfrak m=(3)=\{0,3,6\}=\{9,3,6\}_{\rm vis}.
\]

Además,

\[
\mathfrak m^2=0.
\]

La igualdad anterior es una igualdad de ideales: todo producto de dos elementos de \(\mathfrak m\) se anula módulo nueve. No debe confundirse con el producto cartesiano \(\mathfrak m\times\mathfrak m\), que contiene nueve posiciones de la APP.

\subsection{Secuencia exacta del multiplicador \(3\)}

La multiplicación por tres,

\[
M_3:R\longrightarrow R,
\qquad x\longmapsto3x,
\]

tiene

\[
\ker M_3=\mathfrak m,
\qquad
\operatorname{im}M_3=\mathfrak m.
\]

Por ello induce un isomorfismo de espacios sobre \(\mathbb F_3\),

\[
\overline M_3:R/\mathfrak m\xrightarrow{\sim}\mathfrak m,
\qquad
\overline x\longmapsto3x.
\]

En representantes visuales,

\[
\overline0\longmapsto9,
\qquad
\overline1\longmapsto3,
\qquad
\overline2\longmapsto6.
\]

Así, la banda \(3,6,9\) no es un adorno decimal. Es una copia interna del cociente ternario alojada en el radical nonádico.

\subsection{Órbita cíclica \(369\to693\to936\to369\)}

La fila de la tabla multiplicativa correspondiente a \(3\) es

\[
3,6,9,3,6,9,3,6,9.
\]

El cambio de fase \(j\mapsto j+1\) produce

\[
369\longmapsto693\longmapsto936\longmapsto369.
\]

Estas tres palabras son las tres lecturas cíclicas de la aplicación

\[
\mathbb F_3\longrightarrow\mathfrak m,
\qquad
\bar j\longmapsto3j.
\]

La fila \(6\), al satisfacer \(6\equiv-3\pmod9\), invierte la orientación del mismo ciclo. La marca \(9\) cumple dos funciones tipadas:

\[
9\bmod9=0,
\qquad
\operatorname{dr}_9(9)=9.
\]

El módulo publica la clase nula; la raíz digital positiva conserva la marca de clausura. Esta diferencia de lectores es precisamente la que permite que \(9\) sea a la vez cero residual y cierre visible.

\subsection{Cruz radical de la APP y reducción coordenada a \(\mathbb F_3^2\)}

Sobre

\[
X_9=R^2,
\]

la unión

\[
\mathscr C_{369}
=(\mathfrak m\times R)\cup(R\times\mathfrak m)
\]

es la cruz formada por las filas y columnas \(3,6,9\). Contiene

\[
3\cdot9+3\cdot9-3\cdot3=45
\]

posiciones. Su intersección central \(\mathfrak m\times\mathfrak m\) contiene nueve posiciones. En el observable multiplicativo, la suma de la cruz es

\[
297=216+81=3(72+27).
\]

La reducción coordenada produce

\[
X_9/(\mathfrak m\times\mathfrak m)\cong\mathbb F_3^2
\]

como conjunto cociente tipado por las clases módulo tres. El lift inverso requiere conservar la coordenada dentro de cada fibra de tres elementos.

\subsection{Núcleo \(7227\), descomposición espectral y enlace con la cruz radical}

El bloque central de la APP es

\[
B_c=
\begin{pmatrix}
7&2\\
2&7
\end{pmatrix}
=7I+2R
=9P_++5P_-.
\]

Su lectura por filas produce \(72\mid27\); su lectura por diagonales produce \(77\mid22\). Se verifican

\[
72+27=77+22=99,
\]

\[
72-27=45=\det B_c,
\]

\[
77-22=55=54+1.
\]

La relación con la cruz radical no es nominal:

\[
3\cdot72=216,
\qquad
3\cdot27=81,
\qquad
3\cdot99=297.
\]

El bloque unitario central y la red no unitaria \(3,6,9\) quedan así enlazados por una identidad local–global exacta. La jerarquía de la diferencia suma–producto conserva sus tipos:

\[
4\quad\text{(salto espectral local)},\qquad
2\quad\text{(acoplamiento fuera de la diagonal)},
\]

\[
6\quad\text{(diferencia comprimida módulo tres)},\qquad
54\quad\text{(despliegue bidimensional)}.
\]


\section{Ley potencia--vacancia para las escalas \texorpdfstring{\(9^n\)}{9 elevado a n} y \texorpdfstring{\(10^n\)}{10 elevado a n}}

\subsection{Funciones de longitud decimal y vacancia acumulada}

Para \(n\ge1\), sea

\[
D_9(n)=\#\operatorname{cifras}(9^n)
=\left\lfloor n\log_{10}9\right\rfloor+1.
\]

La frontera decimal \(10^n\) posee \(n+1\) cifras. Definimos la vacancia acumulada de la potencia nonádica por

\[
V_9(n)=(n+1)-D_9(n).
\]

Introducimos los dos números complementarios

\[
\rho_{\rm vac}=\log_{10}9,
\qquad
\delta_{\rm vac}=1-\rho_{\rm vac}
=\log_{10}\frac{10}{9}.
\]

\subsection{Identidad exacta entre longitud decimal y rotación de vacancias}

\textbf{Teorema.} Para todo entero \(n\ge1\),

\[
\boxed{V_9(n)=\left\lceil n\delta_{\rm vac}\right\rceil},
\]

y, por tanto,

\[
\boxed{D_9(n)=n-\left\lceil n\delta_{\rm vac}\right\rceil+1}.
\]

\textbf{Demostración.} Como \(\delta_{\rm vac}\) es irracional, \(n\delta_{\rm vac}\notin\mathbb Z\). Entonces

\[
\begin{aligned}
D_9(n)
&=\lfloor n(1-\delta_{\rm vac})\rfloor+1\\
&=\lfloor n-n\delta_{\rm vac}\rfloor+1\\
&=n-\lceil n\delta_{\rm vac}\rceil+1.
\end{aligned}
\]

Al restar de \(n+1\) se obtiene la primera identidad. \(\square\)

\subsection{Codificación mecánica de las vacancias por las potencias nonádicas}

Los incrementos de la vacancia son

\[
Z_n=V_9(n+1)-V_9(n)
=\lceil(n+1)\delta_{\rm vac}\rceil
-\lceil n\delta_{\rm vac}\rceil
\in\{0,1\}.
\]

Ésta es exactamente la palabra mecánica de pendiente \(\delta_{\rm vac}\) utilizada por el calendario HMT de vacancias. Por tanto, dicho calendario admite una lectura adicional completamente exacta:

\begin{quote}
\(Z_n=1\) si y sólo si, al pasar de \(9^n\) a \(9^{n+1}\), se pierde un nuevo rango decimal respecto de la frontera \(10^n\to10^{n+1}\).
\end{quote}

La potencia no añade una coincidencia a la teoría de vacancias: proporciona su realización natural como contador de pérdida de rango decimal.

\subsection{Ley \(21/22\) como caso del teorema potencia--vacancia}

Puesto que

\[
\delta_{\rm vac}^{-1}
=21.85434532678283256\ldots,
\]

los índices en los que \(Z_n=1\) están separados por \(21\) o \(22\). Los primeros veintiún exponentes satisfacen

\[
D_9(n)=n,\qquad1\le n\le21,
\]

mientras que

\[
D_9(22)=21.
\]

En particular,

\[
9^9=387\,420\,489
\]

posee exactamente nueve cifras. Esta conservación inicial del rango no es una regla indefinida: el calendario de vacancias determina exactamente cuándo deja de cumplirse.

\subsection{Evaluación exacta en \(n=9^9\) y longitud de \(9^{9^9}\)}

Sea

\[
n_*=9^9=387\,420\,489.
\]

Entonces

\[
n_*\delta_{\rm vac}
=17\,727\,389.3684296412564569\ldots,
\]

y, por ello,

\[
\boxed{V_9(n_*)=17\,727\,390}.
\]

La coordenada conservada después de descontar las vacancias es

\[
n_*-V_9(n_*)
=387\,420\,489-17\,727\,390
=369\,693\,099.
\]

Por tanto,

\[
\boxed{\operatorname{ord}_{10}\bigl(9^{9^9}\bigr)=369\,693\,099},
\]

\[
\boxed{D_9(9^9)=369\,693\,100}.
\]

Se obtienen además las congruencias

\[
9^9\equiv0\pmod9,\qquad
V_9(9^9)\equiv0\pmod9,
\]

\[
369\,693\,099\equiv0\pmod9,\qquad
369\,693\,100\equiv1\pmod9.
\]

Las sumas de cifras correspondientes son

\[
s_{10}(9^9)=45,\qquad
s_{10}(V_9(9^9))=36,
\]

\[
s_{10}(369\,693\,099)=54,\qquad
s_{10}(369\,693\,100)=37.
\]

El retorno \(9\to1\) aparece aquí con tipos precisos:

\begin{enumerate}
\item \(9^9\) y la vacancia acumulada pertenecen a la clase nula módulo nueve;
\item su diferencia, el orden decimal, sigue en la clase nula y posee suma de cifras \(54\);
\item el operador universal \(D=\operatorname{ord}+1\) añade la unidad de clausura y publica la clase positiva \(1\).
\end{enumerate}

\subsection{Dominio de validez de la ley potencia--vacancia}

Queda demostrado:

\begin{itemize}
\item el teorema potencia–vacancia para todo \(n\);
\item la identidad con la palabra mecánica de vacancias;
\item la evaluación exacta en \(n=9^9\);
\item el orden y la longitud decimal de \(9^{9^9}\);
\item las congruencias y sumas de cifras anteriores.
\end{itemize}

No se afirma que \(9\) sea la única base con una propiedad análoga en todo dominio, que cada aparición externa de \(369\) proceda de esta torre ni que la torre sustituya la construcción de las nueve fases o el certificado de generación de \(\pi,e,\varphi\).


\section{Jerarquía operatoria, reversibilidad y memoria genealógica}

\subsection{Traslación, dilatación y exponenciación como niveles operatorios}

Las operaciones elementales pueden verse como acciones:

\[
T_a(x)=x+a,\qquad
M_a(x)=ax,\qquad
P_k(x)=x^k.
\]

Sus inversas parciales son, respectivamente, resta, división y extracción de raíces. El paso de una operación a la siguiente modifica la estructura de las fibras.

\subsection{Reversibilidad estratificada en \(\mathbb Z/9\mathbb Z\)}

\textbf{Proposición.} En el anillo nonádico:

\begin{enumerate}
\item toda traslación \(T_a\) es biyectiva;
\item \(M_a\) es biyectiva si y sólo si \(a\in R^\times\);
\item si \(a\in\mathfrak m\), la multiplicación posee fibras no triviales y su imagen queda dentro de \(\mathfrak m\);
\item para \(x\in\mathfrak m\), \(x^k=0\) para todo \(k\ge2\);
\item sobre \(R^\times\cong C_6\), el mapa \(P_k\) es un automorfismo si y sólo si \(\gcd(k,6)=1\).
\end{enumerate}

\textbf{Demostración.} La primera afirmación usa la estructura de grupo aditivo. La segunda es la caracterización de las unidades. La tercera sigue de \(aR\subseteq\mathfrak m\). La cuarta procede de \(\mathfrak m^2=0\). La quinta es la condición para que la multiplicación por \(k\) sea un automorfismo del grupo cíclico de orden seis. \(\square\)

El resultado produce una jerarquía exacta:

\[
\text{suma: reversibilidad global},
\]

\[
\text{producto: reversibilidad restringida a unidades},
\]

\[
\text{potencia: dinámica cíclica en unidades y colapso nilpotente en }\mathfrak m.
\]

Las raíces son funciones sólo cuando la potencia correspondiente es inyectiva en el dominio elegido; en el dominio total son correspondencias de varias ramas o relaciones vacías.

\subsection{Lectura del Topological Phase Kernel sobre la jerarquía operatoria}

Esta estratificación sugiere una definición funcional:

\begin{quote}
El TPK no añade una operación exterior a suma, producto y potencia. Conserva la hoja, el cociente, el acarreo, la orientación, la ruta y la frontera que cada evaluación ordinaria elimina al identificar varias historias con una misma salida.
\end{quote}

La suma y el producto son las dos primeras evaluaciones sobre la APP. La potencia es composición repetida del producto; la raíz selecciona o reconstruye ramas de esa composición. El registro distingue el valor publicado de la historia operatoria que lo produce.

\subsection{Descenso logarítmico entre niveles consecutivos de la jerarquía}

En un dominio positivo,

\[
\log(xy)=\log x+\log y,\qquad
\log(x^k)=k\log x,\qquad
\log(a^x)=x\log a.
\]

El logaritmo convierte producto en suma, potencia en producto y exponencial en dilatación. El defecto aparece cuando el lector posterior discretiza:

\[
\mathcal L_{10}(a^n)
=\lfloor n\log_{10}a\rfloor+1.
\]

La parte entera introduce el calendario de cruces de frontera. Para \(a=9\), ese calendario es precisamente la palabra de vacancias.

\subsection{Conjugación afín de las traslaciones y dilataciones}

Las traslaciones y dilataciones satisfacen

\[
M_aT_bM_a^{-1}=T_{ab}
\]

cuando \(a\) es invertible. Esta identidad expresa que el producto reorganiza las unidades de suma. Cuando \(a\) deja de ser invertible, la conjugación se rompe porque desaparecen ramas. El radical \(3,6,9\) localiza exactamente ese fallo.


\section{Álgebra de caminos: composición multiplicativa, acción aditiva y superposición}

En un álgebra de caminos, la suma agrega alternativas y el producto concatena trayectorias. Para pesos locales \(w(e)\),

\[
K(x,y)
=\sum_{\gamma:x\to y}
\prod_{e\in\gamma}w(e).
\]

Si se fija una rama logarítmica coherente,

\[
\prod_{e\in\gamma}w(e)
=\exp\!\left(\sum_{e\in\gamma}\log w(e)\right),
\]

y, con la acción discreta apropiada,

\[
K(x,y)=\sum_{\gamma:x\to y}e^{-S(\gamma)/\hbar}.
\]

La ecuación reúne las operaciones:

\begin{itemize}
\item el producto compone pasos sucesivos;
\item el logaritmo acumula la acción local;
\item la exponencial reconstruye el peso de la ruta;
\item la suma agrega todas las alternativas.
\end{itemize}

La APP es el soporte mínimo en el que suma y producto actúan sobre las mismas historias. El TPK añade los datos que impiden que la suma sobre caminos se convierta en una suma sobre valores sin genealogía.

En este marco, una cifra es la salida de un lector; un valor es la proyección arquimediana de una sección; un número HMT conserva la familia compatible de caminos que sostiene ese valor; una partícula conserva esa genealogía bajo una realización espectral persistente.


\section{Raíz interna de \(-1\), estructura exponencial y geometrías elíptica, parabólica e hiperbólica}

\subsection{Raíz interna de \(-1\) en la rama finita}

La matriz de incidencia de Paley proporciona un operador \(A_W\) con

\[
A_W^2=-I.
\]

Esto convierte \(\mathbb F_3^6\) en un espacio tridimensional sobre

\[
\mathbb F_9=\mathbb F_3[\iota]/(\iota^2+1).
\]

La raíz de \(-1\) no se adjunta desde fuera como una notación compleja. Emerge como endomorfismo de un soporte ternario ya construido.

\subsection{Realización continua del tipo cuadrático inducido}

En la rama real elíptica, el generador \(J_{\rm e}\) satisface

\[
J_{\rm e}^2=-I,
\]

y

\[
e^{tJ_{\rm e}}=\cos t\,I+\sin t\,J_{\rm e}.
\]

En particular,

\[
e^{\frac\pi2J_{\rm e}}=J_{\rm e},\qquad
e^{\pi J_{\rm e}}=-I,\qquad
e^{2\pi J_{\rm e}}=I.
\]

Así, \(J_{\rm e}\) es el cuarto de giro interno y la ecuación de Euler es su cierre a media vuelta.

\subsection{Soporte operatorio común de \(3\) lecturas geométricas}

Sea

\[
\mathbb D=\{z\in\mathbb C:|z|<1\}.
\]

Este mismo conjunto contiene el diámetro real \((-1,1)\), el interior complejo, la frontera trigonométrica \(S^1\) y la métrica hiperbólica de Poincaré

\[
ds_{\rm P}^2=\frac{4\ell^2|dz|^2}{(1-|z|^2)^2}.
\]

La aplicación

\[
R_{\pi/2}(z)=iz
\]

preserva \(|z|\) y \(|dz|\). Por consiguiente, es simultáneamente:

\begin{itemize}
\item multiplicación compleja por \(\sqrt{-1}\);
\item rotación trigonométrica de \(90^\circ\);
\item isometría euclídea del disco;
\item isometría hiperbólica de Poincaré;
\item automorfismo de su frontera circular.
\end{itemize}

La coincidencia es literal: el mismo mapa actúa sobre el mismo portador; cambia el lector métrico.

\subsection{Geometrías elíptica, parabólica e hiperbólica de la ecuación de Euler extendida}

Los generadores

\[
J_{+1}^2=-I,\qquad
J_0^2=0,\qquad
J_{-1}^2=I
\]

producen

\[
e^{tJ_{+1}}=\cos t\,I+\sin t\,J_{+1},
\]

\[
e^{tJ_0}=I+tJ_0,
\]

\[
e^{tJ_{-1}}=\cosh t\,I+\sinh t\,J_{-1}.
\]

La misma serie exponencial se separa en tres regímenes según el cuadrado del generador: elíptico, parabólico e hiperbólico. La especialización

\[
e^{2\log\varphi\,J_{-1}}=F_{\rm av}^{\,2}
\]

integra la autoescala de Fibonacci. Euler y Fibonacci aparecen como especializaciones de una familia exponencial cuadrática común.

\subsection{Portador común, operador común y familia exponencial cuadrática}

La fusión geométrica se formula en tres niveles:

\begin{enumerate}
\item \textbf{portador común:} el disco unitario contiene el diámetro real, el interior complejo y el borde circular;
\item \textbf{operador común:} \(z\mapsto iz\) es cuarto de giro, multiplicación compleja e isometría hiperbólica;
\item \textbf{familia común:} la ecuación exponencial trítica reúne las ramas elíptica, parabólica e hiperbólica.
\end{enumerate}

El resultado no identifica las métricas euclídea e hiperbólica. Afirma que ambas son realizaciones distintas sobre un mismo soporte y comparten el automorfismo originado por la raíz interna de \(-I\).


\section{Restricción de la rotación interior, holonomía de borde y torsión mínima}

La construcción contiene dos objetos positivos y diferenciados:

\begin{enumerate}
\item el generador interno \(J_{\rm e}\), con \(J_{\rm e}^2=-I\);
\item el invariante global
\end{enumerate}
   \[
   \Delta_4=
   \frac{\pi-e\log\pi}{270}
   \left(1+\frac{\pi}{729}\right),
   \]
   denominado torsión mínima y utilizado antes de la selección de especie.

La tesis autoral añade que la dinámica rotacional interna, al publicarse en el borde, queda absorbida como torsión mínima. La relación no debe escribirse como una igualdad desnuda \(J_{\rm e}=\Delta_4\): los objetos tienen tipos diferentes. Su forma natural requiere

\[
\text{rotación interior}
\xrightarrow{\operatorname{Res}_{\partial}}
\text{holonomía de borde}
\xrightarrow{\operatorname{Def}}
\text{defecto escalar}.
\]

Sea \(q_\partial\) el lector de frontera, \(\nabla_\partial\) la conexión inducida y \(e_\partial\) el marco discreto. La torsión de borde tiene tipo

\[
T_\partial=d_{\nabla_\partial}e_\partial.
\]

La afirmación autoral adquiere forma probatoria mediante un funcional

\[
\Lambda_\partial:
\operatorname{Tors}(\partial X)
\longrightarrow\mathbb R
\]

tal que

\[
\Lambda_\partial(T_\partial)=\Delta_4
\]

y un teorema que haga proceder \(T_\partial\) de la restricción de \(J_{\rm e}\) a través de los operadores de borde del TPK.

La raíz interna de \(-I\), el cuarto de giro, la ecuación de Euler extendida y la fórmula de \(\Delta_4\) quedan demostrados en sus dominios. El paso de la rotación interior a la torsión mínima de borde exige el morfismo \(\operatorname{Res}_{\partial}\) y el funcional \(\Lambda_{\partial}\) declarados; no se identifica por una igualdad escalar.


\section{Triplicidad generacional, conjugación y mezcla sobre el estado enriquecido}

La frase de Isidor Isaac Rabi puede servir como referencia histórica, pero no debe gobernar el título ni constituir una sección autónoma. La cuestión científica es:

\begin{quote}
¿Qué mecanismo interno distingue tres órdenes estables de realización fermiónica y conserva esa distinción bajo color, conjugación y mezcla?
\end{quote}

La respuesta HMT no es «porque TRIT tiene tres valores». La estructura separa

\[
\operatorname{flav}(X)
=\bigl(\sigma_{ud}(X),g(X),\chi_{\rm fina}(X)\bigr),
\]

donde \(\sigma_{ud}\) distingue ramas, \(g\) distingue generación física y \(\chi_{\rm fina}\) conserva orientación y memoria dentro de una rama. En quarks, la órbita de color se añade después del sabor y antes del carácter.

La triplicidad se obtiene como consecuencia de la estructura enriquecida; sus efectos sobre color, conjugación y mezcla se desarrollan a continuación:

\begin{enumerate}
\item color, sabor, generación y masa son lectores distintos del mismo estado enriquecido;
\item la conjugación partícula–antipartícula actúa sobre la orientación sin borrar la clase generacional;
\item CKM y PMNS transportan bases de estados ya tipados; no crean la triplicidad;
\item la clasificación neutrínica debe derivarse de representación, ocupación y mezcla, no de renombrar profundidades \(G0,\ldots,G4\).
\end{enumerate}

Que los neutrinos sean leptones de espín semientero y obediencia fermiónica es una clasificación física conocida. La aportación estructural HMT consiste en obtener internamente por qué el sector neutro ocupa esa representación, por qué admite tres realizaciones físicas y cómo se acopla al operador de mezcla sin leer sus ángulos externos.


\section{Memoria operatoria del TPK: conservación de residuo, cociente, acarreo, hoja, orientación, ruta y frontera bajo proyección}

\subsection{Pérdida de reversibilidad y levantamiento genealógico}

En la APP, distintas historias pueden producir el mismo residuo. El producto pliega más historias que la suma; la potencia pliega más que el producto; la raíz intenta reconstruir ramas ya identificadas. El TPK conserva el dato que hace posible invertir esa pérdida:

\[
\text{residuo},\quad
\text{cociente},\quad
\text{acarreo},\quad
\text{hoja},\quad
\text{orientación},\quad
\text{ruta},\quad
\text{frontera},\quad
\text{memoria}.
\]

La operación algebraica puede caracterizarse no sólo por su valor, sino por el tipo de genealogía que conserva o destruye.

\subsection{Proyección de historias compatibles sobre el continuo arquimediano}

La recta real aparece cuando una familia de cilindros compatibles se contrae bajo la evaluación arquimediana. El valor real es la sombra de la sección; no su genealogía completa. La proyección excepcional conserva otra parte como incidencia y simetría.

La suma sobre caminos expresa el mismo principio:

\[
\text{alternativas}\xrightarrow{\sum}\text{amplitud global},
\qquad
\text{etapas sucesivas}\xrightarrow{\prod}\text{peso de ruta}.
\]

El logaritmo convierte la composición multiplicativa en acción aditiva; la exponencial la devuelve a amplitud. La estructura discreta del continuo puede formularse como conservación genealógica de todas las rutas antes de su proyección sobre \(\mathbb R\).

\subsection{Distinción entre sección numérica y realización persistente de partícula}

Un número y una partícula comparten invariantes ontológicos mínimos: memoria, orientación, persistencia, clausura, firma y posición en una familia de prolongaciones.

El número publica una coordenada arquimediana. La partícula añade realización espectral, ocupación, propagación y estabilidad. El paso HMT \(\to\) MD cambia el lector aplicado a una misma genealogía enriquecida; no cambia arbitrariamente de objeto.


\section{Cierre conjunto de la red radical, la ley potencia--vacancia, la geometría de Euler y la jerarquía operatoria}

\subsection{Teorema de la red radical nonádica}

La red \(3,6,9\) de la APP es la representación visual del radical nilpotente \((3)\subset\mathbb Z/9\mathbb Z\). Sus fases \(369\), \(693\) y \(936\) realizan el isomorfismo \(R/(3)\cong(3)\) inducido por la multiplicación por tres. Las bandas localizan exactamente el fallo de invertibilidad del producto.

\subsection{Teorema potencia--vacancia}

Para todo \(n\ge1\), el número de rangos decimales perdidos por \(9^n\) respecto de \(10^n\) es \(\lceil n\log_{10}(10/9)\rceil\). Sus incrementos forman la palabra mecánica de vacancias. En \(n=9^9\), la coordenada retenida es \(369\,693\,099\) y la longitud resultante es \(369\,693\,100\).

\subsection{Teorema geométrico de la familia exponencial de Euler}

La raíz interna de \(-I\) proporciona un cuarto de giro. En el disco unitario, el mismo mapa \(z\mapsto iz\) es multiplicación compleja, rotación trigonométrica e isometría de Poincaré. La ecuación exponencial trítica prolonga este régimen a los regímenes parabólico e hiperbólico, incorporando en el último la autoescala de Fibonacci.

\subsection{Principio de conservación operatoria entre niveles}

La jerarquía suma–producto–potencia es una jerarquía de fibras y pérdida de reversibilidad. El TPK es el levantamiento que conserva las coordenadas necesarias para distinguir historias identificadas por los lectores ordinarios.


\section{Hipótesis de tipado y controles negativos}

\begin{enumerate}
\item No identificar \(369\,693\,100\) con una palabra TPK sin declarar el lector entre dominios.
\item No confundir el ideal \(\mathfrak m^2=0\) con \(\mathfrak m\times\mathfrak m\).
\item No confundir módulo nueve con raíz digital: \(9\mapsto0\) en el primero y \(9\mapsto9\) en la segunda.
\item No usar la triplicidad de TRIT como demostración suficiente de tres generaciones físicas.
\item No convertir \(J^2=-I\) en \(\Delta_4\) sin el mapa de borde y el funcional de torsión.
\item No presentar \(369\) como prefijo de \(9^{9^9}\): pertenece a su orden y a su longitud decimal.
\item No reducir las diez mil cifras a una prueba de magnitud: certifican truncamientos extensos de un proceso parametrizable cuya tesis es la prolongación indefinida.
\end{enumerate}

\section{Consecuencias conjuntas para la red radical, la memoria operatoria y la torsión de borde}

El hallazgo principal no es que una torre contenga visualmente las cifras \(3,6,9\). La potencia nonádica y el calendario de vacancias son dos presentaciones del mismo defecto logarítmico, mientras la APP localiza \(3,6,9\) como el radical donde la multiplicación pierde invertibilidad. La torre conecta ambos extremos:

\[
\text{radical nonádico}
\longrightarrow
\text{potencia}
\longrightarrow
\text{logaritmo}
\longrightarrow
\text{vacancia}
\longrightarrow
\text{orden decimal}
\longrightarrow
\text{clausura }9\to1.
\]

La ecuación de Euler extendida añade la segunda dirección: la raíz interna de \(-I\) convierte la estructura discreta en dinámica de rotación y reúne, sobre un soporte común, la recta real, el círculo trigonométrico y el disco hiperbólico. El paso ulterior a torsión mínima debe conservar el mapa de borde ya desarrollado en el corpus, no reemplazarlo por una identificación nominal.
